This is the guide I wanted when I first logged into a supercomputer as a PhD researcher. I trained in geo-informatics and remote sensing, which combine geography, geospatial science, programming and mathematics, yet a shared supercomputer was still new to me. It assumes you have never used a command line. By the end, you will have logged in securely, moved files, installed software, built confidence with seven small practice jobs on real data, and run a parameter sweep of many jobs. A final chapter carries these skills to other clusters, the cloud and everyday Linux administration.

The examples use \href{https://ni-hpc.github.io/nihpc-documentation}{Kelvin2}\footnote{\texttt{https://ni-hpc.github.io/nihpc-documentation}}, the cluster run by the Northern Ireland High Performance Computing service (NI-HPC). Most of what you learn carries over to any Linux cluster that uses the Slurm scheduler. \textbf{Kelvin2 only} marks details specific to that system; check your own cluster\textquotesingle s documentation for the equivalents.

\section{How to use it}\label{how-to-use-it}

Prefer one long page? \href{full/}{Read the whole tutorial on one page}.

Read in order if you are new, or jump in if you are not. Blocks of code are labelled with where they run, because typing a command on the wrong computer is the most common beginner\textquotesingle s mistake:

\begin{itemize}
\tightlist
\item
  \texttt{laptop} means a terminal on your own computer.
\item
  \texttt{cluster} means after you have logged in.
\end{itemize}

Lines beginning with \texttt{\#} are comments, which the computer ignores. In this guide, they explain the line beneath them. Anything in \texttt{\textless{}angle\ brackets\textgreater{}} is a placeholder to replace, without keeping the brackets.

\section{Practice files}\label{practice-files}

The \href{https://ahmadammar.github.io/teaching/hpc/examples/hpc-starter-kit.zip}{starter kit}\footnote{\texttt{https://ahmadammar.github.io/teaching/hpc/examples/hpc-starter-kit.zip}} holds small programs and open data for \href{https://ahmadammar.github.io/teaching/hpc/first-jobs/}{the exercises}\footnote{\texttt{https://ahmadammar.github.io/teaching/hpc/first-jobs/}}. Start there if you learn best by doing.

\section{Other good guides}\label{other-good-guides}

This guide is my own, written from setting up on Kelvin2, and it is not an official NI-HPC document. Always treat the official documentation as the authority on your cluster\textquotesingle s rules.

\begin{itemize}
\tightlist
\item
  \href{https://ni-hpc.github.io/nihpc-documentation}{NI-HPC documentation}\footnote{\texttt{https://ni-hpc.github.io/nihpc-documentation}}: the official reference for Kelvin2.
\item
  \href{https://iamdonovan.github.io/teaching/hpc-intro/index.html}{Introduction to high performance computing}\footnote{\texttt{https://iamdonovan.github.io/teaching/hpc-intro/index.html}} by Bob McNabb: a workshop on the shell and on running jobs, on two other university clusters. I recommend it alongside this guide.
\item
  \href{https://swcarpentry.github.io/shell-novice/}{The Unix Shell}\footnote{\texttt{https://swcarpentry.github.io/shell-novice/}} from Software Carpentry, and \href{https://carpentries-incubator.github.io/hpc-intro/}{Introduction to High-Performance Computing}\footnote{\texttt{https://carpentries-incubator.github.io/hpc-intro/}} from HPC Carpentry.
\end{itemize}

Corrections are welcome: contact details are at the foot of the page.

